\chapter{Studio: run the complete physical account}\label{ch:studio-account}

\begin{intuitionbox}{Physical picture: one action, every question in order}
A four-unit delivery is simple enough to describe in one sentence. A defensible account must still identify its state, units, permission, reference, exact field change, actual quantity, and closing identities. This last studio begins with a blank sheet and reconstructs that chain. The aim is to make the short EBU equation fully readable, not to turn it into an automatic social decision.
\end{intuitionbox}

The word ``run'' in this chapter title means working through an account on paper. The examples are illustrative arithmetic, not trajectories from an executed scientific study. A sequence of specified states can teach telescoping without supplying evidence that a controller will produce such a sequence in a real or simulated world.

Keep the main Ridge--Vale fixture. We will first audit one complete lossless transfer, then a short execution, an external-change variant, and a flawed report. Every variation states which data change. This habit is especially valuable near the end of a long calculation, when familiar numbers can encourage the reader to skip the conditions that made them valid.

\section{Prepare the frozen input card}

Write the card with separate columns for role and value:
\begin{center}
\begin{tabular}{P{0.27\linewidth}P{0.27\linewidth}P{0.29\linewidth}}
\toprule
Object & Declared value & Role\\
\midrule
Coordinate order & Ridge, Vale & Physical identity\\
Frozen state & \((18,2)^{\mathsf T}\X\) & Present stock\\
Reference & \((10,10)^{\mathsf T}\X\) & Evaluation definition\\
Scales & \((2,2)^{\mathsf T}\X\) & Evaluation definition\\
Edge & Ridge to Vale & Directed transformation\\
Transfer efficiency & One & Lossless model\\
Finite route cap & \(6\X\) & Permission input\\
Proposed quantity & \(4\X\) & Candidate action\\
Process burden & Zero & Ideal account\\
External change in epoch & None & Baseline scope\\
\bottomrule
\end{tabular}
\end{center}
The route cap is a physical restriction, the scale is part of evaluation, and the proposed quantity is an actor input. Their units and roles differ. We do not derive four from the potential's maximum or from a flux law. It is the given candidate that this account evaluates.

The card's values are illustrative, not clinical recommendations. A real application would need measurement provenance for stock, engineering support for capacity, a justified reference and scale, and a complete declaration of relevant physical effects. The table makes those dependencies visible without pretending to resolve them by arithmetic.

\section{Check the action before evaluating it}

The physical column is \(S_e=(-1,1)^{\mathsf T}\). The proposed withdrawal four is nonnegative, below route cap six, and below source stock eighteen. Under the simple restrictions stated on the card, the candidate is permitted.

Its increment and successor are
\begin{equation}
\delta=S_eq=(-4,4)^{\mathsf T}\X,
\qquad z'=z+\delta=(14,6)^{\mathsf T}\X.
\end{equation}
Both successor stocks are nonnegative. The source withdrawal is four, the destination receipt is four, and the declared loss is zero. The original and successor totals both equal twenty.

If another physical restriction existed, this check would be incomplete. For example, a clinic capacity below six would forbid the endpoint even though the source and route checks passed. Do not infer absence of a constraint from its absence in the arithmetic. The account is complete relative to the declared small world.

\section{Calculate the before and after potentials}

Before action, the standardized deviations are four and minus four. The two local contributions are eight and eight, giving \(V(z)=16\). After action, the deviations are four stock units and minus four stock units, or two and minus two scales. The local contributions are two and two, giving \(V(z')=4\).

The exact signed result is
\begin{equation}
E=V(z)-V(z')-C=16-4-0=12.
\end{equation}
Read it as a sentence: the specified four-unit lossless transfer reduces the declared Gaussian potential by twelve from the frozen state, with no separate process burden in this ideal account. That sentence preserves the baseline and assumptions that the shorter phrase ``twelve EBU'' leaves implicit.

Now check the result independently through the marginal and curvature. The initial marginal is \((2,-2)^{\mathsf T}\X^{-1}\). Its negative dot product with \(\delta\) is sixteen. The curvature correction is four. The finite value is therefore \(16-4=12\), agreeing with the endpoint calculation.

Two algebraic routes reaching twelve are a useful consistency check. They are not two independent experiments, and agreement does not establish that the reference describes a real clinic. The kind of evidence supplied by a check matters as much as the number of checks.

\section{Record proposal, acceptance, and performance separately}

For the complete-match case, the proposal is four, acceptance is four, and illustrative measured performance is four. A receipt can then retain the planned and performed values together. Each label refers to a different stage even though the numbers happen to agree.

Now change only performance to three. The performed increment becomes \((-3,3)^{\mathsf T}\X\), successor \((15,5)^{\mathsf T}\X\), and potential \(6.25\). The performed result is
\begin{equation}
E_{\rm perf}=16-6.25=9.75.
\end{equation}
The proposal did not create a claim to the missing \(2.25\). It remains a planning record. The receipt should explain why acceptance and performance differed if the available evidence can establish that reason. If it cannot, the discrepancy remains unresolved rather than being erased.

The same principle applies when performance is larger than accepted. An exact endpoint value does not excuse an unauthorized quantity. The physical permission record and the value record can each be assessed, and one can fail while the other is arithmetically correct.

\section{Close the three elementary accounts}

For the complete four-unit case, stock closure is
\begin{equation}
18+2=14+6=20\X.
\end{equation}
Transfer closure is \(4=4+0\): withdrawal equals delivery plus declared loss. Field-and-action closure is
\begin{equation}
E+[V(z')-V(z)]+C=12+(4-16)+0=0.
\end{equation}
The final zero includes the field change. It does not say that another person must lose twelve EBU. No payer--payee relation has been derived.

Local-to-global closure follows because only Ridge and Vale change under a separable potential. If an untouched third cell contributed nine to the global potential before and after, the global difference would be \((16+9)-(4+9)=12\). The untouched nine cancels. This does not permit omission of a changed coupling factor in a different potential. Local evaluation must include every affected contribution.

For the short execution, the same identities close using three and \(9.75\): stock remains twenty, and \(9.75+(6.25-16)=0\). A correct account does not require performance to match the plan. It requires the performed transformation to be represented honestly.

\section{One return action tests the sign convention}

From the complete endpoint \((14,6)\), suppose a separately permitted reverse action returns four units to Ridge. It restores \((18,2)\). The reverse field result is \(4-16=-12\). Forward and reverse values add to zero under the fixed potential and ideal zero-burden assumptions.

If a justified process burden of \(0.2\) accompanies each movement, the values become \(11.8\) and \(-12.2\), summing to \(-0.4\). The physical stocks return to their start, but the account retains the two process burdens. Real resource use cannot be cancelled merely because the main stock coordinate completes a cycle.

This example proves a conditional accounting statement for the specified cycle. It does not prove that all actors will avoid cycles or that a proposed balance policy makes cycles impossible. A policy can permit zero-value or affordable negative-value movements. The dynamics of repeated choice remain a distinct research question.

\section{Reconstruct an external event before action}

Consider a variant whose previous recorded state is the reference \(x=(10,10)\), potential zero. An external conservative displacement moves eight from Vale to Ridge, creating the same frozen baseline \(z=(18,2)\), potential sixteen. The external signed potential increment is
\begin{equation}
D_{\rm ext}=V(z)-V(x)=16.
\end{equation}
The actor then performs the familiar four-unit forward transfer, leaving \((14,6)\), potential four, and recording twelve. The actor has reduced part of the displacement; the remaining potential is four. The identity is
\begin{equation}
V(z')+E-D_{\rm ext}=4+12-16=0=V(x).
\end{equation}
The external change did not earn actor credit. It established the baseline against which the later action was evaluated.

One prospective accounting policy would map signed actor attributions into persistent nonnegative capacity balances and record external increments in a separate audit variable \(J\). If the policy validly assigned this single action's twelve to a balance starting at zero, then \(V+B-J=4+12-16=0\). This is a conditional cancellation under declared update rules. It does not establish ownership, permission to spend, fairness, or a recovery theorem. A full policy must supply those missing definitions before an implementation can use the identity operationally.

\section{A smaller account is sometimes the honest result}

Suppose a real receipt has reliable medicine stock observations but no reliable energy record. The medicine field difference may still be computable under its declared model. It should be reported as that resource-specific result with an incomplete physical coverage note. We should not invent an energy burden to make the receipt look comprehensive, and we should not describe the medicine result as a complete ecological evaluation.

Similarly, a measured endpoint does not identify every cause of its change. If external disturbance and actor action overlap without enough information to separate them, the actor-specific interpretation may remain unresolved even when the total endpoint potential is known. The accounting architecture should be able to preserve the measurement without forcing a stronger causal claim.

These are limitations of coverage and identification, not excuses for arbitrary numbers. They tell us what additional observation or modelling would be needed. An account that states an unresolved component clearly can be more useful than a single polished total that conceals it.

\section{Audit a flawed report line by line}

Consider this report: ``The starting state was eighteen and two, and the route could carry six. Two carriers each proposed four, so both proposals were allowed. Each used the original state and obtained twelve EBU. The combined receipt is twenty-four. Only six units actually arrived, but the two accepted proposals justify payment for eight. We clipped Ridge to twelve to restore conservation. Since potential fell, this proves the policy maintains homeostasis.''

First, the route cap is shared, so a joint proposal of eight fails the six-unit cap. Second, two same-baseline four-unit values cannot be added as the joint result. A lossless eight-unit transfer would reach \((10,10)\), with group field reduction sixteen, not twenty-four, if it were physically permitted under a different cap.

Third, actual delivery six does not confirm an eight-unit lossless transfer. The source and destination measurements must be reconciled. If exactly six moved losslessly from the original state, the correct endpoint would be \((12,8)\), potential one, value fifteen. If eight left but six arrived, the present two-cell lossless model is inadequate and the missing two require an explicit physical account.

Fourth, post hoc clipping replaces a physical observation with a preferred number. Conservation can reveal inconsistency, but it does not tell us which observation to invent. Finally, one positive field difference does not demonstrate the long-run behaviour of a policy. The report leaps from an unverified event to an untested dynamic claim.

An audit should record all these defects without quietly repairing the world to make the final claim pass. Different repairs correspond to different events and therefore different evidence. The first step is to determine which physical account is actually supported.

\section{Independent reconstruction exercise}

Set aside the preceding arithmetic and reconstruct the following case. The state is \((17,3)^{\mathsf T}\X\), references ten at each cell, scales two, route cap five, proposed quantity four, accepted quantity four, and performed quantity \(3.5\X\). There is no external change during the epoch and no process burden. Find the complete performed receipt. Then compare its value with the accepted proposal's value and explain the difference.

Begin with the state, not with a remembered quote. Initial deviations are seven and minus seven, so initial potential is \(98/8=12.25\). The accepted proposal would leave \((13,7)\), potential \(18/8=2.25\), giving ten. The performed quantity leaves \((13.5,6.5)\), with deviations \(3.5\) and \(-3.5\), potential \(24.5/8=3.0625\). Its exact value is \(9.1875\).

The performed transfer closes as \(3.5=3.5+0\), and the endpoint total is twenty. The value closure is \(9.1875+(3.0625-12.25)=0\). The quote-to-performance difference is \(10-9.1875=0.8125\). It is the field value of the unperformed half unit from the live short-execution endpoint, not half of the original average value.

Check that last statement independently. At \((13.5,6.5)\), the edge field is \(1.75\). A further half-unit would have exact value \(1.75(0.5)-(0.5)^2/4=0.875-0.0625=0.8125\). The calculation links the local marginal, finite curvature, and receipt difference in one small example.

\section{Final practice questions}

\begin{enumerate}
\item If the independent reconstruction includes a separately justified burden \(C=0.1875\), what net value and closing identity result?
\item If the actual destination were \(6.4\) rather than \(6.5\), while source stock remained \(13.5\), what fails first in the declared world?
\item If a second actor changes a third, disjoint stock pair, what assumptions let its field result be added to this one?
\item If the reference is revised after performance, should the original receipt be silently recalculated under the new definition?
\item Classify the statement ``the two endpoint calculations agree'' and the statement ``the policy keeps clinics supplied for a year.'' What evidence does each require?
\end{enumerate}

With the added burden, net value is nine and closure is \(9+(3.0625-12.25)+0.1875=0\). The physical endpoint is unchanged. In the second question the endpoint total is \(19.9\), so the declared lossless closed-stock account fails before any value interpretation. An explicit loss carrier or corrected observation would be needed.

For disjoint actions, additivity needs complete disjoint affected factors under the declared potential, compatible baselines, and jointly valid physical permission. Merely drawing edges far apart is insufficient if a potential factor or shared constraint couples them. A changed reference should receive its own version and audit treatment; historical receipts retain the definition under which they were made.

Agreement of two calculations is a mathematical or implementation consistency finding, depending on how it is established. Year-long supply is a claim about a specified process, demand, disturbances, and service criteria. It needs the relevant theorem or registered evidence. The first statement cannot substitute for the second.

\section{Write a receipt that another reader can challenge}

A short final report for the base example might read: ``Illustrative lossless stock transfer, Ridge to Vale; baseline eighteen and two; reference ten and ten; scales two and two; quantity proposed, accepted, and performed four; successor fourteen and six; stock total twenty before and after; potential sixteen before and four after; separate process burden zero by idealization; signed field result twelve.''

That sentence is intentionally specific. A reader can challenge the stock observations, the route permission, the choice of reference, the scope of the process idealization, or the arithmetic without confusing them. A claim such as ``the action helped by twelve'' gives fewer handles for a careful review.

The receipt should also say whether its quantities are illustrative, predicted, or observed. In this book they are illustrative. In a deployed system, the recorded observation method and event identity would matter. Exact arithmetic on an illustrative card is educational evidence about the formula's use; it is not a certificate for a physical shipment.

There is a practical advantage to keeping the longer record even when users see a compact display. If a question arises later, the account can be unfolded. A compact interface need not imply a vague scientific record. The book has taken many pages to teach that unfolding because trust should rest on reconstructible meaning rather than on the impressive appearance of a single number.

\section{What the reader can now reconstruct}

The full account begins with physical identities and ends with a carefully scoped interpretation. Between them lie quantities, graph support, permission, a fixed potential, exact finite differences, performance records, and closing identities. You can now rebuild each link rather than treating EBU as an unexplained number at the bottom of a table.

The final chapter returns to the research programme with that foundation in place. Definitions and conditional identities can already be taught and checked. Claims about recovery, sustained homeostasis, institutional usefulness, and fairness still need the additional work appropriate to them. Keeping those questions open is part of the foundation's strength: the account is clear enough for future work to fail honestly as well as to succeed.
